\section{Relation of the two constructions}
Theorem~\ref{thm:ehh} is a $0+1$ dimensional supersymmetric quantum-mechanical construction whose invariant is a Fredholm/Witten index. Theorem~\ref{thm:tetra-susy} is a $3+1$ dimensional supertranslation construction on a discrete null-generated momentum cone. They share a binary grading and a factorization of a positive object as a square or anticommutator, but no functorial identification between the two sectors is proved here.

\section{Nested tetrahedral sections and Euler's Infinity}
One may consider nested tetrahedral sections
\[
\mathcal T_n=r_n\operatorname{conv}\{\mathbf n_1,\ldots,\mathbf n_4\},
\qquad r_n\to0,
\]
around the fixed timelike axis $(1,0,0,0)$. This provides a mathematically clean carrier for an ``Euler Infinity'' programme. However, the scale law $r_n$ must be derived from a declared dynamics. No specific law, including $r_n=e^{-n}$, follows from Euler's identity alone.

The exact supersymmetry theorems above do not require an infinite nest. A canonical inter-section evolution and its compatibility with the spinor/supercharge structures remain \textbf{OPEN / MODEL-LEVEL}.

\section{Proof firewall and conclusion}
The exact results are:
\begin{enumerate}[label=\textbf{E\arabic*.}]
\item the unilateral-shift $\mathcal N=2$ algebra;
\item $\Delta_W(H_k)=\indop(S^k)=-k$ with the declared grading convention;
\item the regular-tetrahedron isotropy and causal-cone identity of Theorem~\ref{thm:cone};
\item rank-one spinor factorization of the four null generators;
\item the $\mathcal N=1$ supertranslation algebra after explicit CAR modes are supplied.
\end{enumerate}

Not established are: physical supersymmetry of Nature; a theory of everything; derivation of the CAR modes from tetrahedral geometry alone; a unique supercharge from a bare pointed complex plane; a canonical scale law for the infinite tetrahedral nest; or any consequence for RH, Collatz, Twin Prime, or another open problem.

The compact operator result is
\[
\boxed{
\text{unilateral shift}+\mathbb Z_2\text{ grading}
\Longrightarrow \mathcal N=2\text{ SUSY},
\qquad
\Delta_W=\indop(S^k)=-k.}
\]
The Hilbert-Hotel defect is therefore literally a supersymmetric index in the declared operator model.

Independently, the tetrahedral spinor/CAR construction gives
\[
\boxed{
\{Q_\alpha,\bar Q_{\dot\beta}\}
=2\sigma^\mu_{\alpha\dot\beta}P_\mu}
\]
on the tetrahedral momentum cone. This upgrades the tetrahedral spinorial picture from analogy to an exact algebraic supersymmetry realization while preserving the crucial boundary that the fermionic degrees of freedom are supplied rather than derived.
